\input format.tex

\usepackage{graphicx}
\graphicspath{{cores/}}

\usepackage{environ}
\usepackage{colortbl,array,booktabs}
\usepackage{tabularx}

\colorlet{TablaBordeSuperior}{topcolor}
\colorlet{TablaBordeInferior}{topcolor}
\colorlet{TablaCentroSuperior}{blue!1}
\colorlet{TablaCentroInferior}{blue!20}
\colorlet{FuenteCabeceraTabla}{white}

\newcolumntype{M}[1]{>{\centering\arraybackslash}m{#1}}
%\newcommand{\tabularxcolumn}[1]{>\arraybackslash}m{#1}}

\tcbset{rtab/.style={
freelance,
frame code={
 \path[top color=topcolor,bottom color=topcolor]
   ([yshift=-#1*(\baselineskip+2pt)]interior.north west) --
   ([yshift=-#1*(\baselineskip+2pt)]interior.north east) {[sharp corners]--
    ([yshift=3pt]interior.north east) --
    ([yshift=3pt]interior.north west)} -- cycle;

  },
interior code={},
 }
}

\newcommand\fuentecabecera[1]{\textcolor{black}{\textbf{#1}}}

\begin{document}

\vspace*{4mm}
%% 各章节
\setlength{\arrayrulewidth}{.2pt}
\fontsize{8.8pt}{11pt}\selectfont
\color{gray2}

\vspace*{3mm}

\noindent{\fontsize{9pt}{11pt}\selectfont\bf{ 二、肠道菌群相对丰度}}

%\vspace*{-4.25mm}
\fontsize{8pt}{11pt}\selectfont
\arrayrulecolor{gray2}
\begin{longtable}{|m{4cm}<{\centering}|m{3cm}<{\centering}|m{3cm}<{\centering}|m{4cm}<{\centering}|}
\hline
\begin{minipage}{4cm}\begin{center}{\vspace*{2mm} {\lantxh\bf 检测内容} \vspace*{2mm}}\end{center} \end{minipage} &
\begin{minipage}{3cm}\begin{center}{\lantxh\bf 检测结果}\end{center} \end{minipage} &
\begin{minipage}{3cm}\begin{center}{\lantxh\bf 检测值}\end{center} \end{minipage} &
\begin{minipage}{4cm}\begin{center}{\lantxh\bf 参考范围}\end{center} \end{minipage} \\
\hline
\begin{minipage}{4cm}\begin{center}{\vspace*{2mm} \lantxh 拟杆菌门(Bacteroidetes) \vspace*{2mm}}\end{center} \end{minipage} &
\begin{minipage}{3cm}\begin{center}{\lantxh 正常}\end{center} \end{minipage} &
\begin{minipage}{3cm}\begin{center}{\lantxh 45.79 {\%}}\end{center} \end{minipage} &
\begin{minipage}{4cm}\begin{center}{\lantxh ≤ 63.44 {\%}}\end{center} \end{minipage} \\
\hline
\begin{minipage}{4cm}\begin{center}{\vspace*{2mm} \lantxh 厚壁菌门(Firmicutes) \vspace*{2mm}}\end{center} \end{minipage} &
\begin{minipage}{3cm}\begin{center}{\lantxh 正常}\end{center} \end{minipage} &
\begin{minipage}{3cm}\begin{center}{\lantxh 37.18 {\%}}\end{center} \end{minipage} &
\begin{minipage}{4cm}\begin{center}{\lantxh ≤ 58.43 {\%}}\end{center} \end{minipage} \\
\hline
\begin{minipage}{4cm}\begin{center}{\vspace*{2mm} \lantxh 变形菌门(Proteobacteria) \vspace*{2mm}}\end{center} \end{minipage} &
\begin{minipage}{3cm}\begin{center}{\lantxh 正常}\end{center} \end{minipage} &
\begin{minipage}{3cm}\begin{center}{\lantxh 10.79 {\%}}\end{center} \end{minipage} &
\begin{minipage}{4cm}\begin{center}{\lantxh ≤ 13.45 {\%}}\end{center} \end{minipage} \\
\hline
\begin{minipage}{4cm}\begin{center}{\vspace*{2mm} \lantxh 蓝藻门(Cyanobacteria) \vspace*{2mm}}\end{center} \end{minipage} &
\begin{minipage}{3cm}\begin{center}{\lantxh 异常}\end{center} \end{minipage} &
\begin{minipage}{3cm}\begin{center}{\lantxh 3.10 {\%}}\end{center} \end{minipage} &
\begin{minipage}{4cm}\begin{center}{\lantxh ≤ 0.60 {\%}}\end{center} \end{minipage} \\
\hline
\begin{minipage}{4cm}\begin{center}{\vspace*{2mm} \lantxh 迷踪菌门(Elusimicrobia) \vspace*{2mm}}\end{center} \end{minipage} &
\begin{minipage}{3cm}\begin{center}{\lantxh 异常}\end{center} \end{minipage} &
\begin{minipage}{3cm}\begin{center}{\lantxh 0.17 {\%}}\end{center} \end{minipage} &
\begin{minipage}{4cm}\begin{center}{\lantxh ≤ 0.00 {\%}}\end{center} \end{minipage} \\
\hline
\begin{minipage}{4cm}\begin{center}{\vspace*{2mm} \lantxh 放线菌门(Actinobacteria) \vspace*{2mm}}\end{center} \end{minipage} &
\begin{minipage}{3cm}\begin{center}{\lantxh 正常}\end{center} \end{minipage} &
\begin{minipage}{3cm}\begin{center}{\lantxh 0.01 {\%}}\end{center} \end{minipage} &
\begin{minipage}{4cm}\begin{center}{\lantxh ≤ 5.58 {\%}}\end{center} \end{minipage} \\
\hline
\begin{minipage}{4cm}\begin{center}{\vspace*{2mm} \lantxh 柔膜菌门(Tenericutes) \vspace*{2mm}}\end{center} \end{minipage} &
\begin{minipage}{3cm}\begin{center}{\lantxh 正常}\end{center} \end{minipage} &
\begin{minipage}{3cm}\begin{center}{\lantxh 0.01 {\%}}\end{center} \end{minipage} &
\begin{minipage}{4cm}\begin{center}{\lantxh ≤ 0.57 {\%}}\end{center} \end{minipage} \\
\hline
\begin{minipage}{4cm}\begin{center}{\vspace*{2mm} \lantxh 厚壁菌门/拟杆菌门(FBR) \vspace*{2mm}}\end{center} \end{minipage} &
\begin{minipage}{3cm}\begin{center}{\lantxh 正常}\end{center} \end{minipage} &
\begin{minipage}{3cm}\begin{center}{\lantxh 0.81}\end{center} \end{minipage} &
\begin{minipage}{4cm}\begin{center}{\lantxh ≤ 4.93}\end{center} \end{minipage} \\
\hline
\end{longtable}

\vspace*{0mm}

\begin{spacing}{1.25}
\noindent\fontsize{9pt}{11pt}\selectfont {人体肠道中超过98{\%}的菌群归属于以下四个菌门，拟杆菌门、厚壁菌门、变形菌门和放线菌门，厚壁菌门/拟杆菌门（FBR）比值根据研究表明与肥胖风险正相关。您的肠道菌群相对丰度性与参考人群趋势一致，说明您的肠道菌群组成结构及稳定性良好，有利于肠道及人体健康。您的肠道菌群不易受到外界不良因素的干扰，但请避免不良作息、刺激性食物、药物滥用等，这样有助于保持良好的肠道菌群相对丰度，维持肠道微生态健康。}

\end{spacing}

\vspace*{8mm}

\noindent{\fontsize{9pt}{11pt}\selectfont\bf{ 肠型分析}}

\vspace*{4mm}

\begin{minipage}[c]{.25\textwidth}
\begin{center}
\noindent\includegraphics[scale=0.5]{enterotype_2.pdf}
\end{center}
\end{minipage}
\begin{minipage}[c]{.47\textwidth}
\noindent\fontsize{9pt}{11pt}\selectfont {经检测，您的肠型为\RNum{2}型，普氏菌型}
\end{minipage}

\vspace*{3mm}
\begin{spacing}{1.25}
\noindent{\ding{84}\ding{84}\ding{84}} \\
\noindent\fontsize{9pt}{11pt}\selectfont {肠型\RNum{1}，肠道中菌属组成以拟杆菌属为主（Bacteriodes），含有合成核黄素（B2）、泛酸（B5）的菌类，合成维生素B7（生物素）与维生素K的酶较多，
主要吸收降解人体无法消化的复杂多糖，产生短链脂肪酸等有益于肠道健康的物质，同时，菌属中少数菌种具有一定致病性（脆弱拟杆菌等）。
研究表明，日常补充较多蛋白质和动物脂肪的人群拟杆菌属含量更高。} \\
\noindent\fontsize{9pt}{11pt}\selectfont {肠型\RNum{2}，肠道中菌属组成以普氏菌属为主（Prevotella），合成维生素B1（硫胺素）与维生素K的酶较多，帮助降解蛋白及部分碳水化合物，
研究表明，素食为主特别是高纤维素饮食习惯的人群普氏菌属含量更高。} \\
\noindent\fontsize{9pt}{11pt}\selectfont {肠型\RNum{3}，肠道中菌属组成以瘤胃球菌属为主（Ruminococcus），瘤胃球菌属主要分解碳水化合物（纤维素），少数也可以分解黏蛋白。
该肠型的人群较少，一般饮食结构比较极端的情况下才可能会出现该肠型。}
\end{spacing}

\end{document}

